% =====================================================================
%  part2intro.tex —— 第二学期补充材料简介（原书 p.130 章首之前，"Short
%  introduction to Part II"，无编号、散文体，两段）
%  说明：本文件位于 \part{第二学期补充材料} 之后、第 9 章之前，只放这一篇
%  非编号导言。原书此页下半即第 9 章标题，其内容见 chunks/ch09.tex。
% =====================================================================
% [原文 p.130]
\chapter*{第二学期补充材料简介}
\addcontentsline{toc}{chapter}{第二学期补充材料简介}
\markboth{第二学期补充材料简介}{}

我大约十年前开始认真关注代数拓扑，主要原因是它能系统地分析量子场论中的量子反常。今天我们已经理解，$d$ 维量子反常是 $d+1$ 维可逆拓扑相的边界现象。用代数拓扑解释这层关系，是我开设这套讲义的初衷；但只有一个学期，材料只够讲到真正讨论这一关系之前，不过关于阿贝尔 Chern--Simons 项量子化的讨论与它关系密切。Part II 是本课程非正式的第二学期，我打算讨论可逆拓扑相与量子反常的代数拓扑处理。

到目前为止，本系列讲义总是先较详细地介绍数学概念，再讨论它们在理论物理中的应用。在 Part II 中，我仍会偶尔先介绍数学概念，但很可能先给出物理动机，再进入数学定义。
